%% kapa: exam=1 round=26 year=2015 period=2 form=A subject=appraisal_law subject_ko=부동산관계법규 seq=29 number=29
%% source: Q-Net 감정평가사 기출문제 (cst00309) artlSeq=5138001 file=extracted/1008/5138001/제26회_감정평가사 1차_기출/제26회_감정평가사 1차_기출/2교시_감평_A형.pdf
%% conversion: kordoc 4.24.51 -> tools/gen1.py (auto); review=proofread-sonnet; flags=none
\begin{question}{29}
건축법령상 건축물의 높이 제한에 관한 설명으로 옳은 것은? (단,「건축법」제73조에 따른 적용 특례 및 조례는 고려하지 않음)
\begin{choices}
  \item 허가권자는 같은 가로구역에서 건축물의 용도 및 형태에 따라 건축물의 높이를 다르게 정하여서는 아니된다.
  \item 가로구역별 건축물의 높이를 지정하는 경우에는 지방건축위원회의 심의를 거치지 아니한다.
  \item 가로구역을 단위로 하여 건축물의 높이를 지정ㆍ공고함에 있어, 건축물의 높이는 지표면으로부터 그 건축물의 상단까지의 높이로 산정한다.
  \item 가로구역별로 건축물의 높이를 지정ㆍ공고할 때에는 해당 가로구역의 상ㆍ하수도 등 간선시설의 수용능력을 고려하여야 한다.
  \item 일반상업지역에서 하나의 대지에 두 동 이상의 공동주택을 건축하는 경우에는 채광의 확보를 위하여 높이가 제한된다.
\end{choices}
\end{question}
\answer{4}
